\documentclass[11pt]{article}
\usepackage[margin=0.8in]{geometry}
\usepackage{amsmath,microtype}
\newcommand{\guideheading}[1]{\par\medskip\noindent{\large\bfseries #1}\par\smallskip}
\begin{document}
\begin{center}
{\Large Guide: An Inverse Step Before Target Convergence}\par\medskip
Collatz proof project\quad---\quad September 5, 2026
\end{center}
\guideheading{What changed}
The earlier inverse excursion began after both base partners reached the
trivial cycle. That could not provide the missing convergence proof for
an unknown target. The new family uses an immediate predecessor of the
source and a short, exact merge. It never assumes the target reaches one.

\guideheading{The induction step}
The transfer question at parameter $u$ asks whether convergence of $3u+2$
implies convergence of $27u+20$. Lean now verifies the reduction
\[
 u=59+192Q\quad\longrightarrow\quad v=39+128Q<u
\]
for every natural quotient $Q$.

The smaller partner at $v$ maps to the smaller partner at $u$ in one step.
Thus knowing the latter converges also proves convergence of the former.
Applying transfer at $v$ supplies its larger partner. That auxiliary larger
partner and the original larger partner then meet: after seven and six
steps respectively, both equal $227+729Q$.

\guideheading{The premise remains explicit}
The proof uses transfer at the strictly smaller parameter $v$. It is an
induction step, not an unconditional proof of transfer on every member of
the family. To obtain a full induction, all parameters still need valid
base cases or decreasing recursive rules.

\guideheading{How the search extends it}
For parameters $u=3w+2$, the immediate predecessor exposes $v=2w+1$.
After applying that smaller transfer, the remaining pair differs by three:
$81w+71$ and $81w+74$.
The search looks for exact merges that preserve both the base endpoint
and its dependence on the quotient.

Through depth twelve, 104 rules cover 407 of the 4096 tested binary
residues of $w$. The other 3689 remain unresolved. These counts concern
one class modulo three and are not global Collatz coverage.

\guideheading{What has been checked}
Lean proves the predecessor identities, the strict parameter decrease,
and the entire displayed family for arbitrary $Q$. Three tests check the
finite search against direct trajectories and replay every saved rule
on several quotient lifts. The full census is not a Lean coverage theorem.

\guideheading{What remains open}
This removes the target-convergence assumption for one useful kind of
certificate. It does not cover the other two parameter classes modulo
three, which include the previously identified $2^k-1$ parameters.
A full proof needs more than extending this one template. Collatz remains
unproved.

\guideheading{Verification files}
\begingroup\raggedright
Build \texttt{Collatz.EarlyInverseBridge} from \texttt{lean/}.
Run \texttt{python3 analysis/early\_inverse\_bridge\_search.py} and
\texttt{python3 -m unittest discover -s analysis -p test\_early\_inverse\_bridge\_search.py}.
\par\endgroup
\end{document}
